% ---------------------------------------------------------------------------
% The algebra of agent reductions, the decision function, and the context bound.
% Included from main.tex after the collectives section.
% ---------------------------------------------------------------------------

\section{The Algebra of Agent Reductions}
\label{sec:algebra}

Naming an operation \code{allreduce} is worth something only if the runtime is
free to choose how to perform it.  MPI bought that freedom with one sentence in
the standard: the operator passed to a reduction must be associative, and the
implementation may therefore evaluate it in any order consistent with
associativity, and with commutativity when the user declares it through
\code{MPI\_Op\_create}'s \code{commute} flag~\cite{mpi41}.  Thirty years of
collective optimisation --- recursive doubling, Rabenseifner's
reduce-scatter/allgather decomposition, ring algorithms, hardware offload ---
all sit inside the licence that sentence grants~\cite{thakur2005collectives,
rabenseifner2004reduction,chan2007}.

The freedom is not free of consequence even in MPI.  Because reordering is
permitted, a floating-point sum is not bitwise reproducible across process
counts, and the standard declines to promise that it will be.  Practitioners
accept this because the error is bounded and well understood.

For agents the same trade appears with the bound removed.  The reductions a
harness actually wants --- merge these partial glossaries, reconcile these two
API designs, synthesise these eight reviews --- are evaluated by a language
model.  Such an operator is at best approximately associative: merging $A$ into
$B$ and then $C$ need not produce what merging $B$ into $C$ and then $A$
produces, and the difference is not a rounding error but a different document.

\ampi{} therefore makes the algebra an explicit, machine-read part of the
operator declaration, and lets it \emph{constrain the algorithm}:

\begin{center}
\begin{tabular}{@{}lll@{}}
\toprule
declared & admissible schedules & depth \\
\midrule
assoc.\ and comm. & any tree, any order, block-decomposed & $O(\log p)$ \\
assoc.\ only      & any tree, rank order                  & $O(\log p)$ \\
neither           & canonical linear order                & $O(p)$ \\
\bottomrule
\end{tabular}
\end{center}

The third row is the interesting one, and it is what MPI never had to write
down because MPI simply refused non-associative operators.  \ampi{} accepts
them, because refusing them would exclude most of the operators a harness
author wants, and pays for them in depth.  Declaring an operator associative
when it is not is the agent analogue of lying to \code{MPI\_Op\_create}: the
program still runs, and the answer silently depends on the agent count.

\subsection{Structural and semantic operators}

Operators come in two kinds, and the distinction is what makes the depth
argument have teeth.

A \textbf{structural} operator executes inside the library.  It costs no
tokens, is exactly reproducible, and its algebra can be established by
inspection.  \ampi{} predefines eleven: \code{AMPI\_CONCAT} (associative, not
commutative), \code{AMPI\_UNION} and \code{AMPI\_BAG} (both), \code{AMPI\_MERGE\_JSON}
(a deep merge whose scalar conflicts resolve to the lexicographically smaller
encoding, which is an arbitrary but total, associative and commutative rule),
the numeric and logical operators, \code{AMPI\_MAXLOC} and \code{AMPI\_MINLOC}
with MPI's lower-rank tie-break, and \code{AMPI\_VOTE}.

A \textbf{semantic} operator cannot execute inside the library, because
evaluating it requires a model.  \ampi{} handles this with an
\emph{upcall}: the library suspends the collective, unwinds to the binding, and
hands the calling agent the two operands together with the operator's
instructions.  The agent evaluates the merge in its own turn, submits the
result with \code{ampi op-submit}, and re-issues the identical collective call,
which resumes at exactly the step where it stopped.  This is the direct
analogue of MPI invoking a user function registered through
\code{MPI\_Op\_create}; the only difference is that the callback runs inside a
model turn and is billed in tokens.

Resumption is what makes the upcall practical, and it is why the collective
layer is implemented as a persisted step machine rather than a straight-line
procedure --- the same design LibNBC uses to make collectives
nonblocking~\cite{hoefler2007nbc}.  Each rank's participation in a collective
compiles to a list of send, receive and combine steps, and the program counter
lives in the device.  One mechanism then buys three properties: collectives
that can be nonblocking, collectives that survive an agent process restarting
mid-operation, and collectives that can suspend for a model upcall without the
other ranks noticing.

\code{AMPI\_VOTE} deserves separate mention because it is the operator with no
MPI counterpart at all.  Pairwise majority is not associative, so \ampi{} does
not implement it pairwise.  Instead the reduction accumulates a commutative
multiset and a \emph{finalize} step takes the mode, reporting the tally and the
agreement fraction rather than only the winner.  This monoid-plus-projection
shape is the honest way to express quorum operators without misdeclaring their
algebra, and it is the agent analogue of algorithm-based fault
tolerance~\cite{huang1984abft}: disagreement among redundant executors is
surfaced, not averaged away.

\subsection{A cost model in tokens, turns, and context}

The Hockney model that underlies MPI collective analysis charges $\alpha$ per
message and $\beta$ per byte, and the analyses that follow from it are about
trading latency terms against bandwidth terms~\cite{hockney1994,chan2007}.  The
substitution for agents is nearly mechanical and one term is new.

Let $p$ be the communicator size and $n$ the payload in tokens.  Write
$\alpha$ for the per-step dispatch latency, $\beta$ for the seconds per token
transferred, and $\gamma$ for the seconds an operator evaluation costs when it
requires a model turn.  Then a schedule of depth $d$, moving $v$ tokens per
rank, with $e$ operator evaluations on the critical path, takes
\begin{equation}
T \;=\; d\,\alpha \;+\; v\,\beta \;+\; e\,\gamma .
\label{eq:cost}
\end{equation}
The fourth quantity, $C$, is not a cost but a \emph{constraint}: the context
window of a single rank.  A schedule whose peak residency
$R$ --- the greatest number of tokens a rank must hold simultaneously ---
exceeds $C$ is inadmissible at any price.

Table~\ref{tab:allreduce-cost} instantiates
Equation~\ref{eq:cost} for the allreduce algorithms \ampi{} implements.  The
columns are the standard ones with bytes replaced by tokens, plus the two that
matter here: $e$, because when $\gamma$ is a model turn it dominates every
other term by three orders of magnitude, and $R$, because it decides
feasibility rather than speed.

\begin{table}[t]
\centering
\small
\caption{Allreduce schedules under the agent cost model. $R$ is peak resident
tokens; $e$ is operator evaluations on the critical path. The last column is
what a semantic operator actually pays, since $\gamma \gg \alpha, \beta$.}
\label{tab:allreduce-cost}
\begin{tabular}{@{}lcccc@{}}
\toprule
schedule & depth $d$ & $v$ & $R$ & $e$ \\
\midrule
linear              & $p$        & $n$              & $2n$   & $p-1$ \\
binomial tree       & $2\lg p$   & $n \lg p$        & $2n$   & $\lg p$ \\
recursive doubling  & $\lg p$    & $2n \lg p$       & $2n$   & $\lg p$ \\
ring                & $2(p{-}1)$ & $2n\frac{p-1}{p}$ & $3n/p$ & $p-1$ \\
Rabenseifner        & $2\lg p$   & $2n\frac{p-1}{p}$ & $2n/p$ & $\lg p$ \\
\bottomrule
\end{tabular}
\end{table}

\subsection{The context bound decides feasibility}

Two rows of Table~\ref{tab:allreduce-cost} differ in a way that has no
counterpart in the HPC reading of the same table.  Recursive doubling and the
binomial tree require every rank to materialise the whole vector at every step,
so $R = \Theta(n)$ independently of $p$.  The reduce-scatter family --- the ring
and Rabenseifner's algorithm --- reduces one block at a time and can discard a
block once it has been forwarded, so $R = \Theta(n/p)$ through the reduction.

In MPI this is a bandwidth argument worth a constant factor, and it is why
Rabenseifner's algorithm wins for long messages~\cite{rabenseifner2004reduction,
thakur2005collectives}.  Under a context bound it is a feasibility argument:

\begin{quote}
For a vector of $n$ tokens over $p$ agents each bounded by a context window of
$C$ tokens, an allreduce whose schedule materialises the whole vector is
admissible only when $n \lesssim C/2$, while a reduce-scatter-based schedule
remains admissible as long as $n \lesssim pC/2$.  The reduce-scatter family is
therefore not merely faster at scale; past $n \approx C$ it is the only family
that can run at all.
\end{quote}

This is a claim about schedules, not about a particular implementation, and
\ampi{} enforces it rather than merely documenting it:
\code{analyse\_allreduce} computes $R$ in closed form and marks a schedule
inadmissible when $R > C$, and \code{select\_algorithm} refuses to return an
inadmissible schedule.  Section~\ref{sec:eval-residency} measures $R$ directly
from the step machine and finds the predicted separation.

\subsection{The decision function}

Open MPI ships a \code{coll tuned} decision file and MPICH a set of control
variables; both select a collective algorithm from message size and
communicator size, tuned per machine~\cite{thakur2005collectives,gabriel2004openmpi}.
\ampi's decision function takes two further inputs that HPC does not have --- the
operator's declared algebra and the rank context bound --- and applies three
admissibility rules in order:

\begin{enumerate}
\item re-association requires a declared associative operator, so a
non-associative operator admits only the canonical linear order;
\item block decomposition additionally requires commutativity, because
reduce-scatter reorders operands within a block --- the same restriction MPICH
places on Rabenseifner's algorithm, for the same reason;
\item the peak residency of the schedule must fit the rank context bound.
\end{enumerate}

Among the survivors it minimises Equation~\ref{eq:cost}.  It also returns its
reasoning: \code{ampi plan} and \code{--explain} report every schedule
considered, its closed-form cost, and, for each rejected one, why.  A harness
author who cannot see why their semantic allreduce took $p$ rounds instead of
$\lg p$ cannot fix it, and the fix is usually a one-word change to an operator
declaration rather than a change to the harness.
